\section{考察とおわりに}
本研究の結果は，仕様を与えることそのものよりも，\textbf{旧仕様から新仕様への対応と，その先に必要となる要素をまとめて与えること}が修正率を左右することを示している．BL-Webは本文に到達しながら45.7\%，A4は対応関係を持ちながら14.3\%にとどまった．修正に必要な要素は仕様上の複数箇所に分散しており，それらを構造的に辿ることが有効であると考えられる．

妥当性への脅威として次の3点が挙げられる．第一に，評価が単一のAPIサービスを対象としており，他のREST APIやドキュメント形式へ一般化しない可能性がある．第二に，修正の成否をマーカーの充足によって判定しているため，マーカーで捉えられない修正の不備を見落とす可能性がある．第三に，先行研究とは生成に用いたモデルが異なるため，手法の差とモデルの差が交絡している．

今後の課題は2点である．第一に，現在は移行表を人手で記述しているため，仕様間の対応付けを自動化する必要がある．第二に，複数の評価者による判定や，実際にコードを実行する検証を導入し，評価の信頼性を高める必要がある．
